% FAC-EN-U0053
% Source: fac_body.tex lines 1889--1924.
% Authority: PDF physical pp. 52--53; printed pp. 247--248 (no. 55 complete).
% English source-first translation; mathematical notation and proof order preserved.

\phantomsection\label{fac:III:2:no:55}
\textbf{55. Sections of $\mathcal{F}(n)$.}

We first prove a lemma on affine varieties, entirely analogous to \hyperref[fac:II:3:no:45:lemma:1]{Lemma~1 of
no.~45}.

\phantomsection\label{fac:III:2:no:55:lemma:1}
\textsc{Lemma 1.}
\textit{Let $V$ be an affine variety, let $Q$ be a regular function on $V$,
and let $V_Q$ be the set of points $x\in V$ such that $Q(x)\neq0$. Let
$\mathcal{F}$ be a coherent algebraic sheaf on $V$, and let $s$ be a section
of $\mathcal{F}$ over $V_Q$. Then, for every sufficiently large $n$, there
exists a section $s'$ of $\mathcal{F}$ over all of $V$ such that
$s'=Q^ns$ over $V_Q$.}

By embedding $V$ in an affine space and extending $\mathcal{F}$ by zero
outside $V$, we reduce to the case in which $V$ is an affine space and hence
is irreducible. By \hyperref[fac:II:3:no:45:thm:2:cor:1]{Corollary~1 to Theorem~2 of no.~45}, there is a surjective
homomorphism
\[
\varphi:\mathcal{O}_V^p\longrightarrow\mathcal{F}.
\]
By \hyperref[fac:II:3:no:42:prop:2]{Proposition~2 of no.~42}, $V_Q$ is an affine open subset, and hence
(\hyperref[fac:II:3:no:44:prop:7:cor:2]{no.~44, Corollary~2 to Proposition~7}) there is a section $\sigma$ of
$\mathcal{O}_V^p$ over $V_Q$ such that $\varphi(\sigma)=s$. We may identify
$\sigma$ with a system of $p$ regular functions on $V_Q$. Applying
\hyperref[fac:II:3:no:43:prop:5]{Proposition~5 of no.~43} to each of these functions, we see that there is a
section $\sigma'$ of $\mathcal{O}_V^p$ over $V$ such that
$\sigma'=Q^n\sigma$ over $V_Q$, provided that $n$ is sufficiently large.
Putting $s'=\varphi(\sigma')$, we obtain a section of $\mathcal{F}$ over $V$
such that $s'=Q^ns$ over $V_Q$.

\phantomsection\label{fac:III:2:no:55:theorem:1}
\textsc{Theorem 1.}
\textit{Let $\mathcal{F}$ be a coherent algebraic sheaf on
$X=\mathbf{P}_r(K)$. There is an integer $n(\mathcal{F})$ such that, for
every $n\geq n(\mathcal{F})$ and every $x\in X$, the
$\mathcal{O}_x$-module $\mathcal{F}(n)_x$ is generated by the elements of
$\Gamma(X,\mathcal{F}(n))$.}

By definition of $\mathcal{F}(n)$, a section $s$ of $\mathcal{F}(n)$ over
$X$ is a system $(s_i)$ of sections of $\mathcal{F}$ over the $U_i$
satisfying the compatibility conditions
\phantomsection\label{fac:III:2:no:55:formula:compatibility}
\[
s_i=(t_j^n/t_i^n)s_j
\qquad\text{on }U_i\cap U_j.
\]
\phantomsection\label{fac:III:2:no:55:def:i-th-component}
We shall call $s_i$ the \textit{$i$-th component of $s$}.

On the other hand, since $U_i$ is isomorphic to $K^r$, there are finitely
many sections $s_i^\alpha$ of $\mathcal{F}$ over $U_i$ that generate
$\mathcal{F}_x$ for every $x\in U_i$ (\hyperref[fac:II:3:no:45:thm:2:cor:1]{no.~45, Corollary~1 to Theorem~2}).
If, for some integer $n$, we can find sections $s^\alpha$ of
$\mathcal{F}(n)$ whose $i$-th components are the $s_i^\alpha$, it is clear
that $\Gamma(X,\mathcal{F}(n))$ generates $\mathcal{F}(n)_x$ for every
$x\in U_i$. \hyperref[fac:III:2:no:55:theorem:1]{Theorem~1} will therefore
follow once we prove the following
lemma.

\phantomsection\label{fac:III:2:no:55:lemma:2}
\textsc{Lemma 2.}
\textit{Let $s_i$ be a section of $\mathcal{F}$ over $U_i$. For every
sufficiently large $n$, there exists a section $s$ of $\mathcal{F}(n)$ whose
$i$-th component is equal to $s_i$.}

Apply \hyperref[fac:III:2:no:55:lemma:1]{Lemma~1} to the affine variety
$V=U_j$, the function $Q=t_i/t_j$, and
the restriction of $s_i$ to $U_i\cap U_j$. This is legitimate, since
$t_i/t_j$ is a regular function on $U_j$ whose zero locus is
$U_j-(U_i\cap U_j)$. It follows that there exist an integer $p$ and a section
$s'_j$ of $\mathcal{F}$ over $U_j$ such that
\phantomsection\label{fac:III:2:no:55:formula:lemma2-local-extension}
\[
s'_j=(t_i^p/t_j^p)s_i
\qquad\text{on }U_i\cap U_j.
\]
For $j=i$, this gives $s'_i=s_i$, so the \hyperref[fac:III:2:no:55:formula:lemma2-local-extension]{preceding formula} may be written
\[
s'_j=(t_i^p/t_j^p)s'_i.
\]

Having defined the $s'_j$ for every index $j$ (with the same exponent $p$),
consider
\[
s'_j-(t_k^p/t_j^p)s'_k.
\]
This is a section of $\mathcal{F}$ over $U_j\cap U_k$ whose restriction to
$U_i\cap U_j\cap U_k$ is zero. Applying \hyperref[fac:II:3:no:43:prop:6]{Proposition~6 of no.~43} to it, we
see that, for every sufficiently large integer $q$,
\phantomsection\label{fac:III:2:no:55:formula:lemma2-compatibility}
\[
(t_i^q/t_j^q)\bigl(s'_j-(t_k^p/t_j^p)s'_k\bigr)=0
\qquad\text{on }U_j\cap U_k.
\]
If we now put $s_j=(t_i^q/t_j^q)s'_j$ and $n=p+q$, the \hyperref[fac:III:2:no:55:formula:lemma2-compatibility]{preceding formula}
becomes
\phantomsection\label{fac:III:2:no:55:formula:gluing}
\[
s_j=(t_k^n/t_j^n)s_k,
\]
and the system $s=(s_j)$ is a section of $\mathcal{F}(n)$ whose $i$-th
component is equal to $s_i$, as required.

\phantomsection\label{fac:III:2:no:55:corollary}
\textsc{Corollary.}
\textit{Every coherent algebraic sheaf $\mathcal{F}$ on
$X=\mathbf{P}_r(K)$ is isomorphic to a quotient sheaf of a sheaf
$\mathcal{O}(n)^p$, for suitable integers $n$ and $p$.}

By \hyperref[fac:III:2:no:55:theorem:1]{the preceding theorem}, there is an
integer $n$ such that
$\mathcal{F}(-n)_x$ is generated by $\Gamma(X,\mathcal{F}(-n))$ for every
$x\in X$. In view of the quasi-compactness of $X$, this is equivalent to
saying that $\mathcal{F}(-n)$ is isomorphic to a quotient sheaf of
$\mathcal{O}^p$, for a suitable integer $p\geq0$. It follows that
\[
\mathcal{F}\simeq\mathcal{F}(-n)(n)
\]
is isomorphic to a quotient sheaf of
$\mathcal{O}(n)^p\simeq\mathcal{O}^p(n)$.
